\documentclass[11pt,a4paper]{article}
\usepackage[utf8]{inputenc}
\usepackage[T1]{fontenc}
\usepackage[ngerman]{babel}

\usepackage{helvet}
\renewcommand{\familydefault}{\sfdefault}

\usepackage{geometry}
\geometry{a4paper, left=2.5cm, right=2.5cm, top=3cm, bottom=3cm}

\usepackage{xcolor}
\usepackage[most]{tcolorbox}
\usepackage{hyperref}
\usepackage{tikz}

\definecolor{uniblue}{RGB}{0, 90, 160}
\definecolor{terminalbg}{RGB}{30, 30, 30}
\definecolor{terminalheader}{RGB}{50, 50, 50}
\definecolor{promptgreen}{RGB}{78, 154, 6}
\definecolor{promptblue}{RGB}{52, 101, 164}
\definecolor{warnred}{RGB}{204, 0, 0}
\definecolor{warnbg}{RGB}{255, 240, 240}


\hypersetup{
	colorlinks=true,
	linkcolor=uniblue,
	urlcolor=uniblue
}

\newcommand{\macbuttons}{%
	\tikz[baseline=-0.6ex]{
		\fill[red!80!white] (0,0) circle (3.5pt);
		\fill[yellow!90!white] (0.4,0) circle (3.5pt);
		\fill[green!80!white] (0.8,0) circle (3.5pt);
	}%
}

\newtcolorbox{terminalbox}[1][]{
	enhanced,
	colback=terminalbg,
	colframe=terminalheader,
	colbacktitle=terminalheader,
	coltext=white,
	fontupper=\ttfamily\small,
	arc=6pt,
	boxrule=1pt,
	toptitle=4pt, bottomtitle=4pt,
	left=12pt, right=12pt, top=10pt, bottom=10pt,
	title={\macbuttons\hspace*{0.5cm}\textcolor{gray!40}{\textsf{\textbf{#1}}}},
	drop shadow={black!30}
}

\newcommand{\prompt}{\textcolor{promptgreen}{student@linux} \textcolor{promptblue}{\char`\~ \$} }

\newtcolorbox{infobox}[1][Wichtiger Hinweis]{
	enhanced,
	colback=warnbg,
	colframe=warnred,
	borderline west={4pt}{0pt}{warnred}, 
	fonttitle=\bfseries\sffamily,
	title={#1},
	coltitle=warnred,
	attach title to upper,
	after title={\\[0.5em]},
	boxrule=0pt,
	arc=3pt,
	left=12pt, right=12pt, top=8pt, bottom=8pt
}

\newcommand{\step}[2]{
	\vspace{0.6cm}
	\noindent\fcolorbox{uniblue}{uniblue}{\textcolor{white}{\textbf{\textsf{ #1 }}}} \textcolor{uniblue}{\textbf{\Large #2}}
	\vspace{0.2cm}\newline
}
\begin{document}
	
	\begin{center}
		{\Huge \textbf{\textsf{\textcolor{uniblue}{CrypTool 2 auf Linux}}}}\\[0.3cm]
		{\large \textsf{Installations- und Startanleitung}}\\[0.4cm]
		\noindent\textcolor{uniblue}{\rule{\textwidth}{1.5pt}}
	\end{center}
	Diese Anleitung führt dich Schritt für Schritt durch den Download, die Installation und den Startprozess von CrypTool 2 unter Linux.
	
	\step{Schritt 1}{ZIP-Datei herunterladen}
	Lade das ZIP-Archiv aus dem Kurs (in Grips) herunter und warte, bis der Download vollständig abgeschlossen ist.\\
	\textbf{Direktlink:} \href{https://grips-link.de}{\texttt{https://grips-link.de}}
	
	\noindent\textcolor{uniblue}{\rule{\textwidth}{1.5pt}}
	
	\step{Schritt 2}{Archiv entpacken}
	In deinem Downloads-Ordner befindet sich nun ein .zip-Archiv. \\
	\textbf{Variante A:} Entpacke die Datei mit einem Doppelklick oder Rechtsklick $\rightarrow$ „Hier entpacken“. Öffne den neuen Cryptool2-Ordner. \\
	\textbf{Variante B:} Wenn du direkt im Terminal arbeitest, kannst du das Archiv mit folgenden Befehlen entpacken und den Ordner betreten:
	
	\begin{terminalbox}[]
		\prompt cd \textasciitilde/Downloads\\
		\prompt unzip Cryptool2.zip -d Cryptool2\\
		\prompt cd Cryptool2
	\end{terminalbox}
	
		\noindent\textcolor{uniblue}{\rule{\textwidth}{1.5pt}}
		
	\step{Schritt 3}{Startdatei ausführen}
	Um die Installation anzustoßen, muss das Skript cryptool.sh ausgeführt werden. \\\textbf{Variante A:} Starte die Datei mit einem Rechtsklick $\rightarrow$ „Auführbare Datei starten“. \\
	\textbf{Variante B:} Direkt im Terminal:
	
	\begin{terminalbox}[]
		\prompt ./cryptool.sh
	\end{terminalbox}
	
	\begin{infobox}[Fehlende Berechtigungen?]
		Falls die Fehlermeldung „Keine Berechtigung“ (Permission denied) auftaucht, fehlen der Datei die nötigen Ausführungsrechte. Führe deshalb im Terminal folgendes aus:
		\vspace{0.2cm}
		\begin{terminalbox}[]
			\prompt chmod +x ./cryptool.sh\\
			\prompt ./cryptool.sh
		\end{terminalbox}
	\end{infobox}

	\noindent\textcolor{uniblue}{\rule{\textwidth}{1.5pt}}

	\step{Schritt 4}{Installation abschließen}
	Verfolge nun die Ausgaben im Terminal. Das Skript richtet Wine, Winetricks und alle benötigten Bibliotheken automatisch ein. 
	
	\begin{itemize}
		\item \textbf{Bestätigungen:} Sollte der Paketmanager nachfragen, bestätige mit y.
		\item \textbf{Passworteingabe:} Da Systempakete installiert werden, fragt das Terminal nach deinem Linux-Benutzerpasswort. Tippe das Passwort einfach ein und drücke Enter.
	\end{itemize}
	Dieser Prozess kann einige Minuten in Anspruch nehmen und auch Fehler schmeißen (einfach ignorieren lol)
	Sobald der Prozess durchgelaufen ist, startet CrypTool 2 automatisch. Sollte dies nicht passieren, führe ./cryptool.sh einfach noch einmal aus.
	
	\noindent\textcolor{uniblue}{\rule{\textwidth}{1.5pt}}
	\vspace{0.2cm}
	
	\textbf{\Large \textcolor{uniblue}{Zukünftiges Starten von CrypTool 2}} \\\\
	Jedes Mal, wenn du CrypTool 2 neu starten möchtest, musst du die cryptool2 einfach nur erneut ausführen. Das System erkennt automatisch, dass bereits alles installiert ist, überspringt das Setup und startet die Software direkt.
	
\end{document}